\chapter{Modelo empírico para branding orientado pela sustentabilidade ambiental da Floresta Amazônica - Uma modelagem de equações estruturais híbridas com abordagem de fsQCA}

\begin{flushright}
Luiz Diego Vidal Santos  \footnote[1]{Programa de Pós-graduação em Ciências da Propriedade Intelectual-PPGPI/Universidade Federal de Sergipe-UFS. *E-mail:  \href{mailto:vidal.center@academico.ufs.br}{vidal.center@academico.ufs.br}}
\\[0.2cm] 
Francisco Sandro Rodrigues Holanda \footnote[2]{Departamento de Engenharia Agronômica/Universidade Federal de Sergipe, São Cristóvão, Brasil. E-mail \href{mailto:fholanda@academico.ufs.br}{fholanda@academico.ufs.br}}
\\[0.2cm] 
Catuxe Varjão de Santana Oliveira \footnote[3]{Departamento de Ciências da Computação/Instituto Federal de Sergipe, Lagarto, Brasil. E-mail \href{mailto:catuxe.oliveira@ifs.edu.br}{catuxe.oliveira@ifs.edu.br}}
\\[0.2cm] 
Pedro Vinícius Bertulino de Menezes \footnote[2]{Departamento de Engenharia Agronômica/Universidade Federal de Sergipe, São Cristóvão, Brasil. E-mail \href{mailto:pedrovinicius18@academico.ufs.br}{pedrovinicius18@academico.ufs.br}}
\\[0.2cm] 
Ana Paula Schervinski Villwock \footnote[2]{Departamento de Engenharia Agronômica/Universidade Federal de Sergipe, São Cristóvão, Brasil. E-mail \href{mailto:ana.agronomia@gmail.com}{ana.agronomia@gmail.com}}
\\[0.2cm] 
Alceu Pedrotti \footnote[2]{Departamento de Engenharia Agronômica/Universidade Federal de Sergipe, São Cristóvão, Brasil. E-mail \href{mailto:alceupedrotti@gmail.com}{alceupedrotti@gmail.com}}
\end{flushright}

% resumo em português
\begin{resumoumacoluna}

Utilizando a teoria do comportamento planejado, é possível estabelecer uma conexão valiosa entre uma organização e seus consumidores, que vai além da aparência e envolve emoções, relações e experiências com a marca. Com base nesses princípios, o objetivo deste trabalho foi apresentar um modelo empírico orientado por evidências a partir da percepção dos empreendedores e consumidores de marcas relacionadas à sustentabilidade ambiental da Amazônia do Brasil. O modelo proposto utiliza modelagem mista, a partir da análise dos dados baseada em análise comparativa qualitativa fuzzy-set (fsQCA), utilizando dois conjuntos de dados originários de pesquisas independentes e obtidas a partir de 8 empresas que operam na floresta amazônica e 535 potenciais consumidores. A partir dos dados observados, foram identificadas as motivações do consumo intrínsecas nos consumidores, as práticas adotadas pelos empreendedores e o papel das comunidades tradicionais e o meio florestal na concepção da marca. Os resultados demostram que para o interesse de compra de marcas sustentáveis ser positivo, a partir da percepção de consumidores a marca deve explorar as informações referentes ao envolvimento com as comunidades tradicionais, percepção da qualidade da marca e inovação da marca, deixando claro também o posicionamento estratégico adotado, uma vez que os construtos se apresentam como centrais da principal solução. Este artigo apresenta informações relevantes para empreendedores de marcas sustentáveis em florestas tropicais, formuladores de políticas empresariais em empreendedorismo sustentável de base comunitária, pesquisadores e interessados em gestão de marca como estratégia para impulsionar pequenos negócios com práticas sustentáveis e envolvimento com comunidades tradicionais. As descobertas e recomendações do estudo podem promover práticas empresariais mais sustentáveis e socialmente responsáveis, fortalecendo a economia local e contribuindo para a preservação ambiental.
 \vspace{\onelineskip}
 
 \noindent
 \textbf{Palavras-chave}: Fuzzy set; Desenvolvimento Sustentável; Manejo florestal; Inovação; Marketing; Brand equity ; Amazônia do Brasil.

\noindent
Artigo sob avaliação no periódico: Quantitative Marketing and Economics
%\{\color{blue}{\underline{https://doi.org/10.53660/CLM-106-127B}}}


%\begin{figure}[H]
 % \centering
  %\caption*{Resumo gráfico 5º capítulo}
  %\includegraphics[scale=0.14]{Imagens/resumo_5.pdf}
  %\fonte{Elaborado pelos autores, 2023}
  %\label{figres:05}
%\end{figure}

%\href{https://empreendasustentavel.com/resumo_5.html}{\color{blue}{\underline{Download: Baixe a imagem em tamanho real}}}

\end{resumoumacoluna}

\section{Introdução}

Definir o crescimento sustentável requer um esforço observacional e prático, pois este ambiente passa atualmente por profundas transformações emergidas a partir de demandas pertinentes à sociedade. O progresso, que antes se apresentava majoritariamente como produção de subsistência e extrativista, hoje dá lugar a um complexo sistema florístico-estrutural e social \cite{Casarejos2018}. Neste sentido, mostra-se importante, ter a compreensão de que definir o desenvolvimento do meio florestal com apenas um conceito seria uma proposição simplista. Buscando contemplar este pressuposto, o desenvolvimento sustentável é muitas vezes visualizado como o conhecido tripé da sustentabilidade (\textit{triple bottom line}), ou seja, ambientalmente correto, socialmente justo e economicamente viável \cite{Nations2015TransformingDevelopment}.

Em seu conjunto, o Triângulo da Sustentabilidade constitui um mecanismo que ajuda a canalizar as lutas, interações e coligações dos movimentos sociais envolvidos na luta pela sustentabilidade, e quando focado na sustentabilidade ambiental busca um processo conjunto entre a manutenção dos ecossistemas naturais e remanescentes (Nascimento 2012). Consonante às ideias propostas pelo tripé da sustentabilidade, estão os Objetivos para o Desenvolvimento Sustentável (ODS) \cite{FAO2014}. Assim como o triângulo da sustentabilidade, os mecanismos funcionais dos ODS são baseados em uma abordagem integrada e multissetorial, envolvendo ampla participação de diferentes atores (setor público ou privado), promovendo sinergias entre a produção agrícola, saúde, setores da educação, do desenvolvimento social e da economia \cite{Kroll2019SustainableSynergies}.

A partir dos dois mecanismos de governança sustentável é possível compreender que, mostra-se crucial adotar uma abordagem integrada e multissetorial para alcançar o desenvolvimento florestal sustentável \cite{Smith2012}. Isso inclui a proteção dos espaços existentes \cite{Billah2022SaltIndicators}, a recuperação de áreas degradadas \cite{Appanah2016ForestForests.}, a adoção de práticas agropecuárias sustentáveis como a silvicultura de base comunitária e da agricultura familiar \cite{Lohmus2016ThePopulations,Bressane2016EnvironmentalSustainability}, que além de manter o ambiente sustentável, são também fundamentais para a segurança alimentar e a economia do país. O desenvolvimento sustentável, como discutido anteriormente, não se limita apenas à proteção ambiental, mas também inclui aspectos sociais e econômicos \cite{Wang2020,Calloway2019EnvironmentalJournal.}, e a percepção do empreendedor quanto à relevância das práticas sustentáveis. No entanto, por motivos econômicos e de capacidade competitiva, esses agentes econômicos são fortemente impactados pelos critérios da economia de mercado aberto \cite{Adzic2008StrategySolutions}. A busca por inovação de processos gerenciais e de marketing pode ser uma solução para alcançar níveis mais altos de competitividade.

O Brasil apresenta diversos biomas de extrema importância para o ecossistema global que merecem ser pautados com tino sustentável, e dentre eles a tem-se Floresta Amazônica. A Floresta Amazônica atualmente é o maior bioma tropical do mundo, abrangendo uma área de aproximadamente 6.7 milhões de quilômetros quadrados em nove países da América do Sul, dos quais o Brasil é o maior detentor, com cerca de 60\% da floresta distribuída por nove estados da federação \cite{Malhi2008ClimateAmazon}. 

Este vasto ecossistema é considerado um patrimônio nacional e global, oferecendo uma infinidade de serviços ambientais, sociais e econômicos. Além da importância para fauna e flora mundial, a floresta atua como um enorme sumidouro de carbono, armazenando aproximadamente 86 bilhões de toneladas de carbono a partir da estocagem de carbono no solo (Ferreira et al. 2018). Além dos valores ecológicos, a Floresta Amazônica tem um grande valor para as comunidades originarias da região \cite{Walker2002LandCycles}. No Brasil a Floresta Amazônica compõe cerca de 350 grupos étnicos e uma população indígena estimada em 1,6 milhão de pessoas na região Amazônica \cite{IBGE2010OIndigena}. Demonstrando que, a Floresta Amazônica é um ecossistema de grande importância para o Brasil, desempenhando funções ecológicas, sociais e econômicas cruciais.


A aplicação do conceito de desenvolvimento sustentável na região amazônica do Brasil não difere quando se considera o tamanho regional e dos enormes dificuldades para implementação. A região amazônica apresenta desafios únicos quanto ao manejo florestal sustentável devido ao seu vasto tamanho, vasta biodiversidade e complexidade socioeconômica da região \cite{Cudney-Valenzuela2023TropicalOpenness,Laurance2002OBrasileira}. A preservação e restauração de áreas florestais degradadas, por exemplo, é um desafio que requer a participação de diferentes setores e atores, incluindo proprietários de terras, governos, organizações não governamentais e a comunidade local \cite{Benayas2009EnhancementMeta-Analysis}. Além disso, o desenvolvimento sustentável na região amazônica requer a consideração dos direitos e necessidades das comunidades originárias e tradicionais que dependem da floresta para sua subsistência \cite{Miranda2020SustainableAmazon,Gouvea2014SustainabilityRegion}. O desafio é encontrar soluções que equilibrem a preservação ambiental com a melhoria da qualidade de vida dessas comunidades.


Neste sentido, para que as pequenas empresas agrícolas e florestais da região amazônica tenham sucesso econômico, considerando as práticas sustentabilidade e o respeito às comunidades ali envolvidas, é necessário apresentarem vantagens competitivas baseadas em inovação, complexidade e/ou sofisticação nos produtos, mudança de mentalidade no manejo da propriedade e dos produtos usados (insumos) e um sistema de compensação dos impactos ambientais que causam \cite{ISI:000579859100002}, ou seja, devem ser inovadoras ou intensivas em conhecimento da economia circular ao mesmo tempo que se mostrem rentáveis economicamente \cite{Barraket2016UnderstandingEnterprises}. 

%Os pequenos negócios que utilizam matéria-prima florestal atualmente buscam associar suas políticas ambientais a estratégias de marketing, por meio do valor de sua marca, baseado na sustentabilidade ambiental . Para isso, a inovação em sistemas gerenciais de marketing para construção de uma consciência de marca é importante para o consumidor considerar a marca entre suas opções de compra . Nos últimos anos, a literatura de marketing mudou para uma lógica que predomina principalmente na percepção do consumidor . 

Pequenas empresas que utilizam matérias-primas florestais buscam atualmente associar suas políticas ambientais a estratégias de marketing baseadas em sustentabilidade ambiental por meio do valor de sua marca \cite{Fardin2020RollerBases}. A literatura sobre marketing verde explora o papel do marketing na promoção do desenvolvimento sustentável \cite{Cheng2022TheAttachment}. O marketing verde, que envolve a promoção de produtos ou serviços com base em seus benefícios ambientais, tem atraído considerável atenção \cite{Parguel2011HowCommunication}. Além disso, tem-se enfatizado o impacto do valor da marca no comportamento do consumidor em relação a produtos sustentáveis \cite{Wang2021TheUK,Hameed2021GreenwashPerspective}.

Para isso, a inovação nos sistemas de gestão de marketing para construir a conscientização da marca é importante, para os consumidores poderem considerar a marca entre suas opções de compra \cite{Cheng2022ImpactSelection}. A literatura de marketing da última década tem visto uma mudança em direção a uma lógica dominante de serviço \cite{Brodie2011CustomerResearch}. 

De acordo com \citeonline{Gupta2016EnhancingRecycling}, essa mudança é refletida em um quadro conceitual baseado na teoria do tripé da sustentabilidade (\textit{triple bottom line}), que assume que a marca pode acelerar a conversão de oportunidades disponíveis em um negócio sustentável, colocando-a à frente de seus concorrentes em desempenho superior. Dessa forma, a marca pode ser um fator estimulante na conquista do desenvolvimento sustentável.

Portanto, associar o \textit{triple bottom line} ao marketing é colocar a sustentabilidade e as pessoas (consumidores e empreendedores) no centro da gestão estratégica da empresa, com base na percepção de que o valor de uma oferta (produto ou serviço) é definido e co-criado a partir do consumidor para os produtos e a compreensão do desempenho da empresa \cite{Becker2020CustomerResearch,Hunt2017StrategicResearch}.

Uma apresentação efetiva das atitudes e políticas de causa ambiental por parte de uma marca pode levar os consumidores a associarem mais facilmente os resultados pretendidos pela marca aos produtos e ao valor justo agregado a eles \cite{Wang2020f,Guzman2017TheFit}. Levando consumidores a gravitarem naturalmente em torno do \textit{brand equity} dessa marca. 

A percepção do consumidor é crucial para criar e associar a imagem da marca a um comportamento planejado, através do conceito de sustentabilidade ambiental. \citeonline{Aaker1992TheEquity} e \citeonline{Keller1993ConceptualizingEquity} concordam que o valor da marca é baseado na consciência, associação, qualidade percebida e lealdade, e pode ser compreendido pela visão do cliente, utilizando estratégias de convencimento \cite{Tasci2018TestingBrands}.


Uma abordagem integrada e multissetorial para o desenvolvimento sustentável da silvicultura, com foco na região amazônica do Brasil, revela-se significativa, uma vez que preenche uma lacuna na literatura ao explorar a interseção entre sustentabilidade ambiental, marketing verde e branding no contexto de pequenos negócios agrícolas e florestais \cite{Chukurna2021MarketingDevelopment}. Isso enfatiza a importância da inovação na gestão da marca e no marketing para alcançar sustentabilidade e competitividade no mercado aberto \cite{Mancuso2021UnderstandingDomain}.

Ao abordar as lacunas mencionadas na literatura, este estudo investiga especificamente e prova ser relevante para uma variedade de partes interessadas. Os empreendedores podem se beneficiar ao entender como integrar a sustentabilidade em suas estratégias de marketing e branding \cite{Oliveira2011Sustentabilidade:Organizacoes}. Os formuladores de políticas podem utilizar os resultados para elaborar políticas que incentivem práticas sustentáveis. Além dos empreendedores, os consumidores se beneficiam ao obter um maior entendimento das práticas sustentáveis das marcas que escolhem apoiar \cite{Minton2012SustainableMedia}.

Neste sentido, a aceitação de produtos sustentáveis pode ser alcançada com técnicas de \textit{Branding}, independentemente da atividade. A autenticidade de marca e estratégias de \textit{Branding} sustentável estão diretamente relacionadas ao desenvolvimento de negócios escaláveis. Neste sentido, esta pesquisa tem por objetivo apresentar um modelo empírico orientado por evidências a partir da percepção dos empreendedores e consumidores de marcas relacionadas à sustentabilidade ambiental da Amazônia do Brasil.

\section{Material e métodos}

\subsection{Desenho da pesquisa e da amostra}
Este estudo foi realizado em 6 (seis) fases. Na Figura \ref{fig_21} é possível compreender o espectro da metodologia de pesquisa adotada neste experimento.

\begin{figure}[H]
  \centering
  \caption{Espectro da metodologia de pesquisa adotada neste experimento.}
  \includegraphics[scale=1.0]{Imagens/fig_21.pdf}
  \fonte{Elaborado pelos autores, 2023}
  \label{fig_21}
\end{figure}

\subsection{Revisão de escopo (Fase I)}

Durante a fase I foi desenvolvida uma revisão de escopo segundo as normas PRISMA, visando identificar os fatores mais importantes referentes às atividades desenvolvidas por negócios sustentáveis relativos à exploração florestal e proteção ambiental, com o potencial de moldar a eficiência de brand equity em marcas locais. A partir dos achados presentes na revisão de escopo foi possível selecionar quatro construtos pertencentes à percepção dos potenciais consumidores frente às marcas analisadas que respondem à intenção de compra (Purchase intention — PI) dos produtos apresentados, sendo eles: posicionamento estratégico de marca (Strategic Brand Positioning — SBP), qualidade percebida da marca (Perceived brand Quality— PBQ), inovação no varejo (Retail Innovation — RI) e Autenticidade da marca (Brand Authenticity — BA).

\subsection{Desenvolvimento da escala (Fase II)}

Na fase II foi desenvolvida e validada a escala psicométrica \textit{brand equity} Sustentável Amazônia (BES-Amazônia) para identificar os fatores mais importantes referentes às atividades desenvolvidas por negócios sustentáveis relativos à exploração florestal e proteção ambiental, com o potencial de moldar a eficiência de \textit{brand equity} em marcas locais. A escala é composta por duas subescalas: escala BES-Amazônia consumidor e a BES-Amazônia empreendedor. 

A subescala BES-Amazônia consumidor foi composta por 22 questões que visam compreender quais construtos influenciam a predileção dos consumidores a marcas sustentáveis e compreender quais os principais atributos e fatores influenciam o potencial da marca a partir da percepção dos consumidores. Ela é composta por quatro construtos: intenção de compra (Purchase intention — PI), posicionamento estratégico de marca (Strategic Brand Positioning — SBP), qualidade percebida da marca (Perceived brand Quality— PBQ), inovação no varejo (Retail Innovation — RI) e Autenticidade da marca (Brand Authenticity — BA). Foi utilizada uma escala do tipo Likert que vai de 1 concordo a 5 concordo totalmente. A validação da escala foi obtida a partir das recomendações propostas no manual padrão para testes educacionais e psicológicos \cite{Association2018}. Para a realizar o teste dessa subescala, foram entrevistados inicialmente 262 consumidores para validação da BES-Amazônia consumidor, tendo com base a percepção das 8 marcas utilizadas na escala BES-Amazônia empreendedor da Amazônia. 

A subescala BES-Amazônia visa compreender a interpretação dos empreendedores sobre as dimensões da sustentabilidade ambiental na região amazônica, com foco nas questões orientadoras das dimensões para o desenvolvimento sustentável. Esta escala baseia-se nas dimensões para o desenvolvimento sustentável propostas por \citeonline{Sachs2002CaminhosSustentavel}\footnote{De acordo com Sachs (2002), o desenvolvimento sustentável só pode ser alcançado quando todas as suas dimensões são consideradas simultaneamente. Essas dimensões incluem a ecológica, que se preocupa com a preservação do meio ambiente e dos recursos naturais; a econômica, que busca o crescimento econômico e a distribuição equitativa de riqueza; a social, que visa a promoção da justiça social e da igualdade; a cultural, que valoriza a diversidade cultural e o respeito às tradições locais; a psicológica, que se preocupa com o bem-estar emocional e mental das pessoas; a territorial, que busca o desenvolvimento equilibrado de todas as regiões; e as políticas nacional e internacional, que promovem a cooperação entre países e a governança democrática em todos os níveis.}. As dimensões econômicas (fonte de matéria-prima, produtos, etc.), social (paridade de gênero, conscientização das comunidades locais, etc.) e ambiental (redução de embalagens, redução das emissões de $CO_2$, etc.) foram avaliadas utilizando uma adaptação do questionário proposto por \citeonline{Postel2006LaFordiste}. As dimensões: culturais, ecológicas, territoriais, políticas (nacional) e políticas (internacional) foram avaliadas utilizando adaptações do questionário proposto por \citeonline{Grigorescu2019KeySustainability}. 

A subescala BES-Amazônia é composta por cinco construtos formativos: Agentes e comunidades envolvidas (Agents and Communities Involved - ACI); Comercialização de produtos e mercado consumidor (Marketing of Prooducts and Consumer Market - MPCM); Disponibilidade de recursos e matéria-prima utilizada (Availability of Resources and Raw Material Used - ARM); Valor adicional da marca (Brand equity - BE) e Visão das atividades florestais sustentáveis (Vision of Sustainable Forestry Activities - VSFA). A subescala BES-Amazônia empreendedor foi testada utilizando 8 (oito) marcas de produtos amazônicos, que produzem considerando a sustentabilidade florestal local.

Foi aplicado o coeficiente de validade de conteúdo (CVC) no instrumento proposto para medir a eficácia da escala com base em avaliações qualitativas de juízes especialistas na área. Para aprimorar o instrumento e contemplar seus objetivos iniciais, essa técnica envolve buscar o consenso de profissionais que irá atribuir valores e opiniões sobre as traduções e adaptações realizadas no instrumento.
A avaliação da escala foi realizada por três juízes especialistas nas áreas interdisciplinarmente correlatas às cegas para garantir a confidencialidade dos participantes e aumentar a fidelidade dos resultados ao final do experimento. Após a validação dos itens, foi realizada a randomização por meio do software Excel. A coleta de dados foi realizada no período de março a maio de 2022. O estudo concentrou-se em 6 cadeias socioprodutivas da região: Turismo, Gastronomia, Biotecnologia, Sistemas Agroflorestais, Produtos Florestais Não Madeireiros e Moda Sustentável.

\subsection{Compreensão da escala (fase III)}

A fase III consistiu na compreensão global da marca utilizando a Escala BES-Amazônia. A compreensão da percepção global da marca a partir do produto foi obtida a partir da submissão do questionário subescala BES-Amazônia consumidores a 535 potenciais consumidores de todas as regiões do Brasil A compreensão da percepção dos empreendedores foi obtida através da subescala BES-Amazônia empreendedores, a partir da entrevista de 8 empreendedores das marcas testadas. Os resultados foram analisados por modelagem por equações estruturais \cite{West2012ModelModeling}, teoria de resposta ao item por método Rasch \cite{Linacre2012WinstepsGuide} e teoria fundamentada \cite{Glaser2017DiscoveryResearch}.

Ao aplicar o modelo Rasch a dados de uma escala Likert, foi possível converter as respostas dos participantes em uma medida linear contínua, que representa o nível de traço latente ($\theta$) subjacente do consumidor, a concordância para com a pergunta \cite{Wright1982RatingAnalysis}. O modelo Rasch pressupõe que a probabilidade de um indivíduo concordar com um item específico depende apenas de sua habilidade latente e da dificuldade do item \cite{Bond2015ApplyingSciences}. Ainda, a utilização deste modelo elimina o pressuposto da tau-equivalência que considera que todos os itens têm a mesma importância para o fator, eliminando os thresholds de cada respondente da pesquisa (consumidores) \cite{Sijtsma2008OnAlpha,Yuan2015}.




\subsubsection{Análise da percepção dos empreendedores das marcas}

As análises da percepção global a partir dos empreendedores, foram conduzidas tendo como base a subescala BES-Amazônia empreendedores por meio de teoria fundamentada. Esta fase estabelece como pressuposto um modelo empiricamente fundamentado que busca compreender como os empreendedores interpretam as dimensões da sustentabilidade ambiental na região amazônica. Especialmente as questões norteadoras das dimensões para o desenvolvimento sustentável descritas por \citeonline{Foley2015UsingExperiences};

A teoria fundamentada se apresenta em quatro etapas: codificação aberta, codificação axial, codificação seletiva alcançando a delimitação da teoria que ao final, se obtêm um guia, contendo códigos associados às respostas dos empreendedores observados \cite{Glaser2017DiscoveryResearchb,Mills2006TheTheory}. As pontuações obtidas em cada código por cada empreendedor entrevistado serviram de base para a fase IV da pesquisa. Estas fases foram implementadas em trabalhos anteriores, sendo este trabalho desenvolvido para contemplar a Fase IV e subsequentes.


\subsection{Concordância dos guias (fase IV)}


Nesta fase, implementada a partir desta pesquisa, as entrevistas dos oito (8) empreendedores entrevistados foram submetidas ao escrutínio de dez (10) especialistas, buscando assim atribuir valores de inclusão em cada cluster presente no guia. 

A pontuação para a concordância seguiu os padrões propostos pela escala Likert com cinco pontos: 1 Não se aplica; 2 Usualmente não se aplica; 3 Ocasionalmente se aplica; 4 Usualmente se aplica; 5 Sempre se aplica. Para avaliar a concordância das pontuações atribuídas pelos especialistas seguindo o guia, foram empregados testes Kappa Simples. A análise do Kappa Simples foi realizada utilizando-se o software SPSS (versão 25.0).

A interpretação dos resultados dos testes Kappa baseou-se nos critérios estabelecidos por \citeonline{Landis1977TheDatab}, os quais estipulam que valores de concordância oscilando entre 0 e 0.2 são considerados insatisfatórios ou fracos; entre 0,21 e 0,4, são tidos como discretos; entre 0.41 e 0.6, caracterizam-se como regulares ou moderados; entre 0,61 e 0,8, são apontados como bons ou substanciais; e, finalmente, aqueles que variam de 0.81 a 1.0 são classificados como ótimos ou excelentes.


\subsection{Análise por fsQCA (fase V)}

Este estudo optou pela análise fsQCA como método agregador dos dados observados, haja vista que ele permite que as variáveis de resultado e preditoras existam em uma escala difusa (contínua) e não apenas em uma escala dicotômica (binária), independente de sua condição inicial (qualitativa ou quantitativa). Neste sentido, o objetivo do fsQCA é transformar os escores brutos das variáveis variadas em escores de associação difusos que variam de 0.0 a 1.0 \cite{Wu2019IdentifyingApproachb}.

Assim, análise fsQCA se mostra como uma nova contribuição para esta linha de pesquisa, uma vez que segue o paradigma da teoria da configuração de conjuntos, que permite analisar interações holísticas entre elementos de natureza confusa e não linear. Neste sentido, os oito construtos (variáveis constituintes) da escala BES-Amazônia foram considerados variáveis de condição para análise de comparação.

Atendendo padrões procedimentais para o método fsQCA, será realizada a calibração para os valores fuzzy. A métrica é calibrada em conjuntos fuzzy com valores que variam de 0 a 1.0 onde 1.0 = pertence ao conjunto completo, 0.5 = pontos cruzados e 0 = nenhuma relação clara. Como todas as variáveis da escala BES-Amazônia estão distribuídas seguindo a escala Likert de 5 pontos, o redimensionamento foi necessário \cite{Chang2014HowSatisfaction}. 

A calibração dos dados resultantes das entrevistas com os potenciais consumidores será realizada tendo como base os valores de traço latente dos respondentes $\theta$ para cada construto a partir do somatório dos itens pertencente. Os valores: Para os construtos os valores pertencentes aos conjuntos completos (1.0) foram considerados os mais altos traços latentes observados ($\theta$ SBP = 5.790; $\theta$ PBQ = 3.720; $\theta$ RI = 7.17; $\theta$PI = 3.49).  Os valores de pontos cruzados (0.5) foram selecionados por mediana ($\theta$ SBP = 1.53; $\theta$ PBQ = 0.22; $\theta$ RI = 3.290; PI $\theta$ = -1.260) e por fim o valor para nenhuma associação (0) foi considerado a partir do menor valor de traço latente obtido em cada construto ($\theta$ SBP = -4.61; $\theta$ PBQ = -6.140; $\theta$ RI = -6.940; PI $\theta$ = -5.77). As calibrações e análises foram realizadas utilizando o software fsQCA v4.0 \cite{Ragin2023AnalyticResearch}. Tendo processado as calibrações das medidas de condição e resultado em pontuações fuzzy, os resultados do fsQCA foram divididos em três partes: análises de necessidade, explicação de configurações e análises de suficiência.

Para a subescala BES-Amazônia empreendedores, foram considerados os valores atribuídos pelos avaliadores especialistas para os índices presentes no guia: cinco (5) = conformidade total, três (3) = interseção e um (1) = não conformidade, seguindo sugestões de \citeonline{Kaya2020AntecedentsfsQCA}, já que todas as medidas apresentaram intervalos aceitáveis, sugerindo uma distribuição normal. 

Na busca por resultados em \textit{fuzzy-set} QCA, para a construção da tabela da verdade foi adotada a abordagem proposta por \citeonline{Rihoux2006QualitativeMethods}, que consiste na identificação de condições centrais presentes em soluções parcimoniosas com consistência acima 0.75. Sendo as condições periféricas eliminadas. A presença de uma condição é indicada por círculos pretos ($\bullet$), enquanto a ausência é representada por círculos riscados ($\bigotimes$) \cite{Ragin2008NetDemonstration}, o não preenchimento indica indiferença do construto na solução. Já distinção entre núcleo e periférico é feita usando círculos grandes e pequenos, respectivamente.

Para seleção dos casos, foi utilizado o valor de cobertura por “Redução Proporcional na Inconsistência” PRI. Para ser aceito na solução foi estipulado o limite mínimo de 0.7 para o escore de consistência do PRI \cite{Greckhamer2018StudyingResearch}. 

\subsection{Proposição do modelo empírico (fase VI)}

Por fim, na fase VI foi formulado um modelo empírico permitindo que pesquisadores e empreendedores construam suas próprias marcas ou auxiliem novas pesquisas multidisciplinares para suas marcas, auxiliando assim essas empresas no desenvolvimento de marcas positivas sem prescindir da responsabilidade ambiental necessária ao desenvolvimento local. O modelo aqui apresentado visa avançar o conhecimento das práticas experimentais de marcas na academia, responder a pesquisas emergentes e vincular este trabalho a ferramentas de gestão ambiental.

Em termos da multidisciplinaridade exigida na área estudada, o modelo propõe uma melhor interação entre pesquisadores e especialistas da área de marketing, \textit{Branding} e profissionais da área de estudos da sustentabilidade. Nesse sentido, o objetivo é fomentar a discussão sobre o tema, identificar os pontos fortes e as limitações das abordagens atuais e propor novas soluções para o desenvolvimento florestal sustentável.

Assim como afirma \citeonline{Mandran2018THEDRE:Model}, sempre que um modelo empírico é desenvolvido, ele deve ser colocado num paradigma epistemológico. Assim, o modelo se assenta num paradigma epistemológico construtivista pragmático que visa esclarecer os seus pressupostos epistemológicos, procurando gerar conhecimento científico a partir de suposições hipotéticas baseadas em experiências a partir de uma realidade vivida \cite{Dewey2009Constructivism}.

A Figura \ref{fig_22} apresenta a rede nomológica de relacionamentos conceituais que fundamentam o modelo de \textit{Branding} sustentável. Foi adotado o S-O-R estabelecido ou a estrutura de entrada-resposta-saída (input–response–output), conforme encontrado em muitos modelos de intenção de compra \cite{Rose2012OnlineOutcomes}. O modelo é composto por três tipos de variáveis: variáveis antecedentes, variáveis componentes do \textit{Branding} e variáveis de desfecho. 

\newpage
\begin{landscape}
\begin{figure}[H]
  \centering
  \caption{Rede nomológica de relacionamentos conceituais que fundamentam o modelo de Branding sustentável}
  \includegraphics[scale=0.3]{Imagens/fig_22.pdf}
  \fonte{\textbf{Elaborado pelos autores} (2023)}
  \label{fig_22}
\end{figure}
\end{landscape}

Foram consideradas variáveis antecedentes às indagações presentes nas escalas, apoiadas na literatura. Os construtos foram desenvolvidos a partir das variáveis antecedentes, consideradas as variáveis formativas, assim como de componentes cognitivos e afetivos da percepção global da marca a partir do produto e a percepção global da marca a partir dos empreendedores. Isso por sua vez levam ao comportamento de desfecho (percepção global da marca): Intenção de compra. A seguir, as hipóteses consideradas:

Intenção de compra. Apresentamos a seguir uma discussão do modelo com hipóteses.

$H_1$ — O posicionamento estratégico de marca a partir de um conjunto baseado na teoria do comportamento planejado influencia positivamente a intenção de compra;

$H_1$ — A qualidade percebida da marca a partir de um conjunto baseado na teoria do comportamento planejado influencia positivamente a intenção de compra;

$H_3$ — A Percepção da Inovação no varejo da marca a partir de um conjunto baseado na teoria do comportamento planejado influencia positivamente a intenção de compra;

$H_4$ — A Percepção da autenticidade da marca a partir de um conjunto baseado na teoria do comportamento planejado é influenciada positivamente pela intenção de compra;

$H_6$ — A ligação entre as marcas e aos agentes e comunidades envolvidas na floresta amazônica a partir de um conjunto baseado na teoria do comportamento planejado influencia positivamente a intenção de compra;

$H_7$ — O conhecimento do consumidor sobre o meio de comercialização e o mercado atuante de marcas sustentáveis da floresta amazônica influencia positivamente a intenção de compra;

$H_8$ — O conhecimento do consumidor sobre a origem dos recursos e matéria-prima utilizada por marcas sustentáveis da floresta amazônica influencia positivamente a intenção de compra;

$H_9$ — O valor adicional atribuído a produtos de marcas sustentáveis da floresta amazônica influencia positivamente a intenção de compra;

$H_{10}$ — A visão dos empreendedores sobre a importância das atividades sustentáveis desenvolvidas pela sua marca influencia positivamente a intenção de compra;

$H_{11}$ — A teoria do comportamento planejado a partir de um conjunto difuso representativo da percepção do consumidor quanto à responsabilidade ambiental do empreendedor a influência a intenção de compra de marcas sustentáveis.

Por se tratar de estudos com seres humanos, a pesquisa foi encaminhada ao Comitê de Ética em pesquisa, aprovado para a aplicação. Os participantes que aceitaram a participação no estudo foram esclarecidos quanto aos objetivos, a coleta de dados, bem como, sobre as informações contidas no Termo de Consentimento Livre e Esclarecido (TCLE). Quanto aos que não aceitaram participar, seus dados foram excluídos da pesquisa.


\section{Resultados e discussão}

\subsection{Validação do instrumento}

A amostra utilizada para o teste da subescala BES-Amazônia empreendedor apresentou uma distribuição por idade desequilibrada. Dos participantes, 42\% estavam entre 18 e 24 anos, 34\% entre 25 e 34 anos, 15\% entre 35 e 44 anos, e 9\% entre 45 e 64 anos. Além disso, a maioria dos participantes (54,5\%) foi do sexo feminino. Quanto à procedência, consumidores de todas as regiões do Brasil estiveram presentes no estudo, sendo 6\% da região norte, 40\% da região Nordeste, 32\% da região Sudeste e 16\% da região Sul do Brasil. Em termos de resposta ao questionário, 104 participantes responderam presencialmente, enquanto 158 responderam de forma online.

A subescala BES-Amazônia consumidores apresentou índices de ajuste da escala foram adequados (c2 = 86.460. gl = 132; p < 0.001; RMSEA = 0.027; BCa 95\% [0.00 — 0.037]; CFI = 0.998; BCa 95\% [0.998 – 1,002]; TLI = 0.997; BCa 95\% [0.997 - 1,002]). Na raiz do erro quadrático médio de aproximação (RMSEA), um valor inferior a 0.08 para o instrumento proposto indica um bom ajuste estatístico para o modelo \cite{West2012ModelModelingb,Tucker1973AAnalysisb}. Índices de ajuste comparativo (CFI) e índice Tucker-Lewis (TLI) acima de 0.95 também indicam um bom ajuste do instrumento, corroborando com o bom desempenho do instrumento.

Quando observada a análise individual da escala, os valores de Bartlett (2926,2, \textit{gl} = 210; p < 0.001) e KMO (0.89) sugeriram a interpretabilidade da matriz de correlação dos itens. Para estabilidade, o alfa de Cronbach apresentou valores aceitáveis (Cronbach's = 0. 925). Tendo a análise paralela sugerindo cinco fatores refletindo o esperado para os construtos teoricamente propostos. 

Ainda sobre a validade do instrumento, o modelo foi submetido a validade convergente que analisa grau de correlação entre o teste e outros testes que medem o mesmo construto. O instrumento apresentou valores aceitáveis para todas as dimensões analisadas SBP (0.845), PBQ (0.964), PI (0.889), BA (0.675) e RI (0.677), já que, a análise convergente deve apresentar valores acima de 0.50 \cite{Maroco2014}.

Os resultados alcançados indicam que a subescala BES-Amazônia apresenta solidez psicométrica, com valores de validação das variáveis dentro do esperado, quando associados ao interesse de potenciais consumidores em relação às marcas analisadas. Dessa forma, a subescala validada demonstrou ser aplicável neste experimento. O Apêndice 1 descreve os itens de medição.

Em relação à validação da subescala BES-Amazônia empreendedor, dois itens que não atingiram valores superiores ou iguais a 0.8 foram agrupados e realocados para correção em pelo menos dois dos critérios estabelecidos (clareza, relevância e coerência). Após as correções, o instrumento apresentou valores adequados de Coeficiente de Validade de Conteúdo (CVC) para utilização prática. O CVC médio da escala completa, considerando os três critérios - clareza, coerência e relevância - foi de 0.88. Além disso, apresentou bons índices de clareza de linguagem (0.78), relevância prática (0.87) e relevância teórica (0.93), resultando em um CVC geral de 0.91 e um erro PeJ de 0.002. Estes resultados demonstram que o instrumento psicométrico BAS-Amazônia a partir das duas subescalas mostrou-se adequado ao público brasileiro. O Apêndice 2 descreve os itens de medição e o Apêndice 3 apresenta o guia final utilizado.

\subsection{MEE da percepção dos consumidores}

Com base nos resultados da Modelagem por Equações Estruturais (MEE) obtida a partir da aplicação da sub escala BES-Amazônia em 543 potenciais consumidores, foi estabelecido o modelo de definição do \textit{brand equity} baseado em marcas relacionadas à sustentabilidade ambiental da Amazônia que influenciam a intenção de compra PI dos consumidores ao produto. 

O modelo estrutural (Figura 3) demostra que a predileção dos consumidores aos produtos de marcas sustentáveis é influenciada direta e positivamente por todas as variáveis estudadas. Tendo os construtos BA: $\beta$ = 0.886; erro padrão (SE) = 0.041; (p < 0.00) e RI: β = 0.881; erro padrão (SE) = 0.074; (p < 0.00) os maiores valores de influência. Os resultados mostraram efeitos significativos, assim como, os caminhos padronizados entre o fator e os indicadores acima de 0.40. portanto, sustentável \cite{Wang2019StructuralModels}.

\newpage
\begin{landscape}
\begin{figure}[H]
\centering
\caption{Contribuição das variáveis posicionamento estratégico da marca (SBP), percepção da qualidade dos produtos (PBQ), Inovação no varejo (RI) e autenticidade da marca (BA) na intenção de compra dos produtos oriundos de marcas sustentáveis da Amazônia (PI). ($\chi$2/gl=4.94; CFI=0.910; TLI=0.981; SRMR =0.781; RMSEA=0.084 (0.071 — 0.109)}
\includegraphics[scale=0.5]{Imagens/fig_13.pdf}
  \label{fig_23}
\end{figure}
\end{landscape}

Quanto a análise de concordância dos juízes acerca do enquadramento dos empreendedores e suas marcas nos índices contidos no guia proposto, a comparação por Kappa de Fleiss para o diagnóstico principal mostrou haver uma confiabilidade moderada entre os três avaliadores (K = 0.645; [IC95\%: 0.524 - 0.710]; z = 9.784; p< 0.001), denotando uma concordância estatisticamente significativa entre os avaliadores. Os resultados da análise de Kappa para categorias individuais mostraram-se também aceitáveis e estatisticamente significativos para todas as classes presentes no guia. 

Os indicadores apresentados nas soluções principais mostraram que a percepção do consumidor a partir do produto apresenta suficiência para o interesse a compra PI, tendo como condições centrais o posicionamento estratégico de marca (SBP), a qualidade percebida da marca (PBQ), a inovação no varejo (RI), porém não necessárias isoladamente para ao interesse à compra. 

Este comportamento pode ser observado no gráfico  dos construtos associados do interesse a compra (Figura 4). A Figura 4 revela que essa relação é de fato assimétrica, por exemplo, a Figura \ref{fig_24}a indica que valores elevados da percepção do consumidor frente à sustentabilidade ambiental da marca levam a valores elevados de interesse a compra (WTP); no entanto, baixos valores da percepção também levam a valores altos e baixos de desempenho da empresa. Portanto, altos valores de percepção da sustentabilidade pelos consumidores são suficientes para o alto desempenho de interesse da marca, mas não são exatamente necessários, uma vez que baixos valores da percepção do consumidor também podem levar a um interesse a aquisição. 

Ao considerar o construto ACI para a solução possível (Figura \ref{fig_24}b) os valores de PI se comporta mais simétrico que os dados anteriores. Este resultado demonstra a forte influência da percepção dos consumidores frente à preocupação dos empreendedores para com o envolvimento de sua marca com as comunidades envolvidas (ACI). 


\begin{figure}[H]
\centering
\caption{(a) Gráfico de interesse a compra PI diante da percepção global da marca a partir do produto e (b) Gráfico XY de interesse a compra PI diante da percepção global da marca sustentável.}
\includegraphics[scale=0.08]{Imagens/fig_24.pdf}
  \label{fig_24}
\end{figure}


A compreensão de cada construto pode ser obtida a através da análise das configurações das soluções possíveis. Nas configurações presentes na Tabela \ref{tab:16} são apresentados os diferentes caminhos necessários e suficientes presentes na interação entre os construtos analisados e a intenção de compra (\textit{Purchase intention} — PI) dos produtos oferecidos pelas marcas sustentáveis. Mais de uma combinação de condições é suficiente para alcançar o resultado. Diferentes combinações mostram uma consistência global satisfatória (em torno de 0. 888) e uma cobertura global também satisfatória (0.892).

Cada coluna representa uma combinação alternativa de condições que se associam ao desfecho intenção de compra (PI). Com base na análise realizada, foi possível identificar cinco soluções alternativas com valor de consistência aceitável, resultando na aceitação dos produtos pelos clientes avaliados. Cada uma dessas soluções representa um agrupamento de construtos que compartilham configurações aceitáveis e de elementos ou antecedentes, os quais estão associados a altos níveis de conformidade, ou seja, concordância total.


\begin{table}[H]
\centering
\caption{Configurações das condições presentes na interação entre os construtos analisados e a PI disposição}
\label{tab:16}
\begin{tabular}{lcccccl}
\hline
\multicolumn{1}{c}{\multirow{2}{*}{Configuração}} & \multicolumn{6}{c}{Solução}                                                                                           \\ \cline{2-7} 
\multicolumn{1}{c}{}                              & I                    & II                   & III                  & IV                   & \multicolumn{2}{c}{V}     \\ \hline
SBP                                               & $\bullet$                    & $\bullet$                    & $\bullet$                    & $\bullet$                    & \multicolumn{2}{c}{$\bullet$}     \\
RI                                                & $\bullet$                    & $\bullet$                    & $\bullet$                    & $\bullet$                    & \multicolumn{2}{c}{$\bullet$}     \\
BA                                                & $\bullet$                    & $\bullet$                    &                      & $\bullet$                    & \multicolumn{2}{c}{$\bigotimes$}     \\
PBQ                                               & $\bullet$                    & $\bullet$                    &                      & $\bullet$                    & \multicolumn{2}{c}{$\bullet$}     \\
ACI                                               & $\bullet$                    & $\bullet$                    & $\bullet$                    & $\bullet$                    & \multicolumn{2}{c}{$\bullet$}     \\
MPCM                                              & $\bullet$                    & $\bigotimes$                    & $\bigotimes$                    &                      & \multicolumn{2}{c}{$\bigotimes$}     \\
ARM                                               & $\bullet$                    &                      & $\bullet$                    &                      & \multicolumn{2}{c}{$\bigotimes$}     \\
BE                                                & $\bullet$                    & $\bullet$                    & $\bullet$                    &                      & \multicolumn{2}{c}{$\bullet$}     \\
VSFA                                              & $\bullet$                    & $\bullet$                    & $\bullet$                    &                      & \multicolumn{2}{c}{$\bullet$}     \\ \hline
Nº   casos (\% cumulativa) &
  \multicolumn{1}{l}{51   (21\%)} &
  \multicolumn{1}{l}{48   (42\%)} &
  \multicolumn{1}{l}{25   (52\%)} &
  \multicolumn{1}{l}{18   (60\%)} &
  \multicolumn{2}{l}{14   (66\%)} \\
Consistência                                      & 0.988                & 0.973                & 0.923                & 0.964                & \multicolumn{2}{c}{0.896} \\
Cobertura bruta                                   & 0.737                & 0.548                & 0.595                & 0.576                & \multicolumn{2}{c}{0.287} \\
Cobertura única                                   & 0.486                & 0.   123             & 0.0253               & 0.011                & \multicolumn{2}{c}{0.010} \\
                                                  & \multicolumn{1}{l}{} & \multicolumn{1}{l}{} & \multicolumn{1}{l}{} & \multicolumn{1}{l}{} & \multicolumn{2}{l}{}      \\ \hline
Consistência   geral                              & 0.888                & \multicolumn{1}{l}{} & \multicolumn{1}{l}{} & \multicolumn{1}{l}{} & \multicolumn{2}{l}{}      \\ \cline{1-2}
Cobertura   geral                                 & 0.892                & \multicolumn{1}{l}{} & \multicolumn{1}{l}{} & \multicolumn{1}{l}{} & \multicolumn{2}{l}{}      \\ \hline
\end{tabular}
\begin{minipage}{\linewidth}
    \raggedright
    $\bullet$ Presença de uma construção | $\bigotimes$ Ausência de uma construção
\end{minipage}
\end{table}


A primeira solução (I) representa o conjunto de maior cobertura (0.737) e consistência (0.988). É possível compreender que, para ocorrerem intenções de compra das marcas sustentáveis com maior assertividade, todos os construtos devem estar presentes, já que essa solução reflete a combinação da presença de todos os construtos analisados. 

A segunda e terceira solução (II e III) reflete a principalmente a negação do construto MPCM, associada a ausência dos construtos ARM (II). Esta solução explica que altas intenções de compra podem ser alcançadas com alta personalização da mensagem, independentemente do que os benefícios relacionados com a inatividade do produto. Outra ocorrência de solução que corrobora com este comportamento é a solução III. 

Esta solução (III) demonstra que a aceitação de produtos sustentáveis pode ser alcançada mesmo com a ausência da percepção pelos consumidores referente a autenticidade da marca (BA) e a qualidade percebida no produto (BPQ). Neste sentido, independentemente da inexistência da percepção de qualidade, da autenticidade e a baixa relação com mercado local, os consumidores, quando bem apresentados ao brand equity, tendem a consumir tais produtos. Resultados similares foram observados por \citeonline{Pappas2021Fuzzy-setMarketing}. 

Esta solução demostra que, mesmo com a ausência dos fatores pertinentes a percepção dos empreendedores, e os consumidores, a partir das informações apresentadas sobre as marcas, estão dispostos a adquirir tais produtos.  \citeonline{Liu2020b} observaram em seu estudo que o conhecimento sobre as atitudes sustentáveis das marcas, quando percebidas pelos clientes (por meio do reconhecimento e imagem da marca), afeta positivamente a percepção da marca nos níveis cognitivo e afetivo. 

Em uma análise comparativa dos resultados obtidos neste estudo com pesquisas anteriores, foram identificadas semelhanças e diferenças significativas. Consistente com os achados de \citeonline{Mohr2001DoBehavior}, este estudo também reconheceu a percepção da qualidade da marca como um fator influente na motivação do consumo. Contudo, uma distinção importante reside no foco da pesquisa. Enquanto \citeonline{Mohr2001DoBehavior} direcionaram sua atenção para a responsabilidade social corporativa e seu impacto no comportamento de compra, o presente estudo concentrou-se especificamente em marcas associadas à sustentabilidade ambiental na região da Amazônia brasileira. Esta particularidade sugere que, embora a qualidade da marca seja um elemento de importância em todos os setores, sua associação com a sustentabilidade ambiental pode exercer um impacto particularmente acentuado.

Assim como as soluções IV e V demonstram que, as percepções da Inovação no varejo (RI), Comercialização de produtos e mercado consumidor (MPCM) e Disponibilidade de recursos e matéria-prima utilizada (ARM) desempenham um papel menor na intenção de compra (BE E VSFA), independentemente da inexistência da percepção global da marca a partir dos empreendedores. A partir destes resultados, é possível rejeitar as hipóteses alternativas H7 e H8. Esta mesma solução demostra ainda que as condições centrais são fatores suficientemente forte, já que, mesmo com ausência da RI, MPCM e ARM ainda há cobertura aceitável. Este resultado pode estar atrelado também a capacidade das mensagens presente no discurso da marca. 

Devido à capacidade reduzida de redação publicitária que está presente em uma embalagem ou logomarca, as marcas tendem a não priorizar a apresentação das práticas sustentáveis e os procedimentos adotadas em torno do gerenciamento de \textit{Branding}, deixando estes conteúdos informacionais disponíveis em seu body copys. Estas dificuldades podem ser mitigadas a partir de construções de narrativas que visam o apelo a sustentabilidade, narrativas estas baseadas em metodologia AINDA \cite{Radu2021EuropeanAccessing}, e a partir desta impulsionar o discurso da marca sustentável sobre a relevância das práticas adotadas pelas empresas, além do engajamento da marca sobre a plena sustentabilidade \cite{Kasuma2022SustainabilityAsia}.  

É digno de nota que assim como observado na modelagem por equações estruturais dos três construtos SBP, RI e PBQ, dois (SBP e RI) mostraram ser condições centrais a existência do PI para as analisadas a partir da análise QCA. Este resultado indica que o conjunto de construtos representativo da percepção do comportamento dos empreendedores sobre a sustentabilidade ambiental e o conhecimento da floresta pouco influencia a intenção de compra de potenciais consumidores.

A qualidade percebida da marca (PBQ), mesmo presente em quatro das cinco soluções concebidas na análise fsQCA, apresenta-se como uma condição periférica para o resultado positivo ao interesse da compra de produtos das marcas sustentáveis. A qualidade do produto se refere a quão bem um produto satisfaz as necessidades do cliente, serve ao seu propósito e atende aos padrões esperados pelos consumidores. Ela é um fator determinante que contribui para a construção da percepção do cliente em relação a produtos de uma marca \cite{Dangi2021CustomerDiscounts}.

Ao avaliar a qualidade do produto a partir de um exercício imaginário, os consumidores consideram vários fatores-chave, incluindo se um produto resolve o seu problema, se ele funciona de maneira eficiente ou atende aos propósitos de satisfação com o seu consumo \cite{Rodrigues2021FactorsConsumers}.

Porém, à medida que a qualidade já se faz presente no produto a partir de acordos sociais (confiança da marca, reputação, atendimento a requisitos sanitários) os consumidores passam a considerar outros fatores relevantes para a escolha como: confiança do consumidor e mudanças no gosto e valoração das práticas éticas das empesas, e inovabilidade do produto e expectativa de compra, estes fatores emergem influenciando fortemente o interesse na compra \cite{Sasmita2015YoungEquity,Nam2011BrandSatisfaction,Mustafa2022EclaircissementAnalysis}.

Este comportamento pode ser explicado a partir da teoria da necessidade de Maslow. Conforme a teoria, as necessidades humanas são organizadas em uma hierarquia, começando com as necessidades fisiológicas básicas e subindo até a autorrealização \cite{Niemela2014MaslowsNeeds}. A percepção da qualidade do produto pode ser um fator importante, porém não decisivo isoladamente para os consumidores quando a estes restam a completude da autorrealização social e pessoal.

Neste sentido, nossas descobertas contribuem para a literatura existente sobre \textit{Branding} e sustentabilidade ambiental, proporcionando uma nova camada de compreensão. Embora estudos anteriores, como o de \citeonline{Sander2021429}, tenham examinado a importância da sustentabilidade para a percepção da marca, nosso estudo é um dos primeiros a explorar essa questão no contexto específico da Floresta Amazônica. Além disso, a aplicação de modelagem mista e análise fsQCA em nossa pesquisa oferece uma abordagem metodológica inovadora para o estudo dessas questões.

\subsection{Modelo empírico}

A Figura 5 apresenta o modelo proposto a partir da caracterização dos principais a partir das evidências observadas nesta pesquisa. A partir dos achados foi possível propor um modelo conceitual baseado na percepção das marcas sustentáveis da Amazônia. Dos cinco construtos formativos da percepção global da marca a partir do produto, três são condições centrais (SBP, PBQ e RI) tornando o construto BA não necessário para a construção do \textit{Branding} para marcas sustentáveis.

As seis primeiras hipóteses indicam que todas as variáveis consequentes formativas da percepção global da marca a partir das informações presentes nos produtos possuem cobertura e consistência apropriadas que fornecem validação empírica parcial das hipóteses levantadas. Seis das onze hipóteses são apoiadas, implicando uma sólida fundamentação teórica da especificação do modelo. O poder de cobertura de cada variável consequente varia de 0.56 a 0.78 (Figura 5). 

\newpage
\begin{landscape}
\begin{figure}[H]
\centering
\caption{Modelo empírico orientado pelo conceito de sustentabilidade ambiental da Floresta Amazônica}
\includegraphics[scale=0.3]{Imagens/fig_25.pdf}
  \label{fig_25}
\end{figure}
\end{landscape}


Os resultados observados sugerem uma relação direta ao “tripé da sustentabilidade” (triple bottom line) a partir da importância do envolvimento das marcas tanto com os fatores ambientais quanto os sociais a partir da percepção de consumidores. Também para ambos os modelos, só é possível a contemplação da sustentabilidade em sua totalidade a partir da conscientização da relevância dos recursos ambientais pelos empreendedores.

Assim, a proposta do modelo se baseia em três pontos principais: 

\begin{enumerate}[label=(\roman*)]
- Uma marca sustentável para alcançar bons resultados quanto ao interesse à compra dos seus produtos deverá em seu \textit{Branding} demostrar sua inovação sustentável ambientalmente aplicada nos produtos;
- O \textit{Branding} de marcas sustentáveis deve sempre demonstrar com clareza a qualidade atribuída aos seus produtos, em todos os meios de contato com o consumidor;
- A marca sustentável presente na Floresta Amazônica deve buscar alinhar-se aos povos originários da floresta, buscando com estes meios e estratégias se manter sustentável ambientalmente não ferindo os conhecimentos e tradições destes povos, à medida que mantêm uma relação precípua e saudável;
\end{enumerate}

\section{Conclusões}

A valoração do produto a partir do \textit{brand equity}, recai na possibilidade de melhor atuação frente aos desafios de manter-se sustentável, uma vez que o movimento em direção à sustentabilidade exigirá mudanças em muitos níveis, então os métodos de produção devem ser mais respeitosos ética e ambientalmente, e os padrões de propagandas e estratégias de divulgação também devem ser alterados.

As hipóteses aceitas neste estudo apontam para a importância das estratégias de \textit{Branding} e da qualidade percebida da marca como fatores que influenciam de maneira positiva a disposição dos consumidores à compra dos produtos apresentados. Esses elementos são fundamentais para construir uma \textit{brand equity} positivo a partir da teoria do comportamento planejado.

A Análise QCA realizada mostrou que a presença dos construtos referentes ao posicionamento estratégico de marca (SBP), inovação no varejo (RI) e envolvimento com as comunidades locais (ACI) são aspectos centrais que contribuem para a intenção de compra de produtos sustentáveis. Estes fatores são cruciais para o sucesso das marcas sustentáveis no mercado. Embora a percepção da qualidade do produto (PBQ) seja vista como uma condição periférica no estudo, ela ainda possui um papel relevante na intenção de compra do consumidor. A garantia e a promoção da qualidade dos produtos são essenciais para estabelecer e manter a confiança do público no portfólio de produtos oferecidos.

O modelo empírico proposto neste estudo estabelece uma relação entre os construtos analisados e a intenção de compra, baseando-se no conceito da hierarquia de necessidades de Maslow e no \textit{triple bottom line}. Este modelo enfatiza aspectos relacionados com as necessidades humanas, o compromisso ambiental, social e econômico das empresas e como tais aspectos afetam a percepção de valor do consumidor em relação à aquisição dos produtos ofertados.

Com base nos resultados obtidos, fica evidente a necessidade de as marcas sustentáveis direcionarem seus esforços para a inovação sustentável nos produtos, a clara comunicação da qualidade dos produtos e o alinhamento com os povos originários da Floresta Amazônica. Ao se dedicar a esses fatores, as empresas estarão em uma posição mais favorável para alcançar bons resultados no que diz respeito ao interesse de compra por parte dos consumidores, fortalecendo assim sua imagem e relevância no mercado.
A contribuição acadêmica desta pesquisa se destaca pelo fortalecimento da sustentabilidade ambiental como uma medida formativa do \textit{brand equity} de marcas que produzem e comercializam produtos florestais de forma sustentável. A marca, a partir desta natureza produtiva, pode usufruir dos benefícios ali presentes, estimulando o engajamento e aceitação dos produtos a partir da percepção dos consumidores. Dessa forma, as práticas sustentáveis se tornam uma base para a criação de novas marcas ambientalmente responsáveis.